\documentclass[11pt,a4paper]{article}

\usepackage[utf8]{inputenc}
\usepackage[T1]{fontenc}
\usepackage[spanish]{babel}
\usepackage{amsmath,amssymb,mathtools}
\usepackage{geometry}
\usepackage{lmodern}
\usepackage{microtype}

\geometry{margin=2.4cm}
\setlength{\parindent}{0pt}
\setlength{\parskip}{0.45em}

\newcommand{\R}{\mathbb{R}}
\newcommand{\sen}{\operatorname{sen}}

\begin{document}

\begin{center}
{\LARGE\bfseries Geometría Lineal}\\[0.4em]
{\Large Tema 2: Geometría analítica del plano y del espacio}\\[0.4em]
{\large Soluciones}
\end{center}

\section{Ejercicio 1}
\subsection{Solución}

La recta pasa por:

\[
A=(1,2,1)
\]

y tiene vector director:

\[
v=(4,5,-1).
\]

Su ecuación vectorial es:

\[
\boxed{
(x,y,z)
=
(1,2,1)
+
\lambda(4,5,-1)
}
\]

Las ecuaciones paramétricas son:

\[
\boxed{
\begin{cases}
x=1+4\lambda,\\
y=2+5\lambda,\\
z=1-\lambda.
\end{cases}
}
\]

La ecuación continua es:

\[
\boxed{
\frac{x-1}{4}
=
\frac{y-2}{5}
=
\frac{z-1}{-1}
}
\]

\section{Ejercicio 2}
\subsection{Solución}

Sean:

\[
A=(1,0,1),
\qquad
B=(0,1,1).
\]

Un vector director es:

\[
\overrightarrow{AB}
=
(-1,1,0).
\]

Por tanto, la ecuación vectorial es:

\[
\boxed{
(x,y,z)
=
(1,0,1)
+
\lambda(-1,1,0)
}
\]

Las ecuaciones paramétricas son:

\[
\boxed{
\begin{cases}
x=1-\lambda,\\
y=\lambda,\\
z=1.
\end{cases}
}
\]

La ecuación continua puede escribirse como:

\[
\boxed{
\frac{x-1}{-1}
=
\frac{y}{1},
\qquad
z=1
}
\]

Como todos los puntos de la recta verifican:

\[
z=1,
\]

la recta es paralela al plano:

\[
z=0.
\]

Por tanto:

\[
\boxed{
r\cap\{z=0\}
=
\varnothing
}
\]

\section{Ejercicio 3}
\subsection{Solución}

El plano pasa por:

\[
A=(1,1,1)
\]

y contiene los vectores:

\[
v=(1,-1,1),
\qquad
w=(2,3,-1).
\]

Las ecuaciones paramétricas son:

\[
\boxed{
\begin{cases}
x=1+\lambda+2\mu,\\
y=1-\lambda+3\mu,\\
z=1+\lambda-\mu.
\end{cases}
}
\]

Un vector normal es:

\[
n=v\times w
=
(-2,3,5).
\]

Por tanto:

\[
-2(x-1)+3(y-1)+5(z-1)=0.
\]

Así:

\[
\boxed{
-2x+3y+5z-6=0
}
\]

\section{Ejercicio 4}
\subsection{Solución}

Sean:

\[
A=(-1,2,3),
\qquad
B=(3,1,4).
\]

Entonces:

\[
\overrightarrow{AB}
=
(4,-1,1).
\]

El plano contiene además:

\[
v=(0,0,1).
\]

Por tanto:

\[
\boxed{
\begin{cases}
x=-1+4\lambda,\\
y=2-\lambda,\\
z=3+\lambda+\mu.
\end{cases}
}
\]

Un vector normal es:

\[
n
=
\overrightarrow{AB}\times v
=
(-1,-4,0).
\]

Podemos tomar:

\[
n=(1,4,0).
\]

Por tanto:

\[
(x+1)+4(y-2)=0.
\]

Así:

\[
\boxed{
x+4y-7=0
}
\]

\section{Ejercicio 5}
\subsection{Solución}

Sean:

\[
A=(1,0,1),
\quad
B=(0,0,2),
\quad
C=(3,2,1).
\]

Se tiene:

\[
\overrightarrow{AB}=(-1,0,1),
\]

\[
\overrightarrow{AC}=(2,2,0).
\]

Un vector normal al plano \(ABC\) es:

\[
\overrightarrow{AB}\times\overrightarrow{AC}
=
(-2,2,-2).
\]

Por tanto, el plano es:

\[
\boxed{
x-y+z-2=0
}
\]

Intersecamos con:

\[
z=3.
\]

Entonces:

\[
x-y+1=0.
\]

Tomando:

\[
y=\lambda,
\]

se obtiene:

\[
x=\lambda-1.
\]

Por tanto, la recta de intersección es:

\[
\boxed{
(x,y,z)
=
(-1,0,3)
+
\lambda(1,1,0)
}
\]

\section{Ejercicio 6}
\subsection{Solución}

Las ecuaciones paramétricas son:

\[
\begin{cases}
x=1+\lambda+\mu,\\
y=2-\lambda+\mu,\\
z=4\lambda-3\mu.
\end{cases}
\]

Los vectores directores son:

\[
v=(1,-1,4),
\qquad
w=(1,1,-3).
\]

Un vector normal es:

\[
v\times w
=
(-1,7,2).
\]

El plano pasa por:

\[
P=(1,2,0).
\]

Así:

\[
-(x-1)+7(y-2)+2z=0.
\]

Por tanto:

\[
\boxed{
-x+7y+2z-13=0
}
\]

\section{Ejercicio 7}
\subsection{Solución}

La recta es:

\[
\frac{x-1}{2}
=
\frac{y+3}{1}
=
\frac{z}{-1}.
\]

Por tanto, pasa por:

\[
P=(1,-3,0)
\]

y tiene vector director:

\[
v=(2,1,-1).
\]

Además, el plano debe pasar por:

\[
A=(2,0,1).
\]

Entonces:

\[
\overrightarrow{PA}
=
(1,3,1).
\]

Un vector normal es:

\[
n
=
v\times\overrightarrow{PA}
=
(4,-3,5).
\]

Así:

\[
4(x-2)-3y+5(z-1)=0.
\]

Por tanto:

\[
\boxed{
4x-3y+5z-13=0
}
\]

\section{Ejercicio 8}
\subsection{Solución}

La recta está dada como intersección de:

\[
2x+3y-z=1
\]

y:

\[
x-y+3z=2.
\]

Los vectores normales son:

\[
n_1=(2,3,-1),
\qquad
n_2=(1,-1,3).
\]

Por tanto, un vector director de la recta es:

\[
v
=
n_1\times n_2
=
(8,-7,-5).
\]

El plano buscado es perpendicular a la recta, por lo que \(v\) es un vector normal al plano.

Como pasa por:

\[
A=(1,3,2),
\]

se tiene:

\[
8(x-1)-7(y-3)-5(z-2)=0.
\]

Por tanto:

\[
\boxed{
8x-7y-5z+23=0
}
\]

\section{Ejercicio 9}
\subsection{Solución}

Tomamos:

\[
A=(0,5,3),
\quad
B=(0,6,4),
\quad
C=(2,4,2),
\quad
D=(2,3,1).
\]

Se tiene:

\[
\overrightarrow{AB}=(0,1,1),
\]

\[
\overrightarrow{AC}=(2,-1,-1),
\]

\[
\overrightarrow{AD}=(2,-2,-2).
\]

Calculamos:

\[
\det
\begin{pmatrix}
0&2&2\\
1&-1&-2\\
1&-1&-2
\end{pmatrix}
=
0.
\]

Por tanto:

\[
\boxed{
A,B,C,D
\text{ son coplanarios}
}
\]

\section{Ejercicio 10}
\subsection{Solución}

La recta \(r\) es:

\[
(x,y,z)
=
(2,5,1)
+
\lambda(-1,1,3).
\]

Sustituyendo en los planos que definen \(s\):

\[
x+y-2z-1=0,
\]

\[
3x+3y-4z-9=0,
\]

se obtiene:

\[
\lambda=\frac23.
\]

Por tanto, ambas rectas se cortan en:

\[
\boxed{
P=
\left(
\frac43,
\frac{17}{3},
3
\right)
}
\]

Un vector director de \(s\) es:

\[
v_s
=
(1,1,-2)\times(3,3,-4)
=
(2,-2,0).
\]

Podemos tomar:

\[
v_s=(1,-1,0).
\]

El vector director de \(r\) es:

\[
v_r=(-1,1,3).
\]

Entonces:

\[
\cos\theta
=
\frac{|v_r\cdot v_s|}
{\|v_r\|\|v_s\|}
=
\frac{2}{\sqrt{22}}.
\]

Por tanto:

\[
\boxed{
\theta
=
\arccos\left(
\frac{2}{\sqrt{22}}
\right)
}
\]

\section{Ejercicio 11}
\subsection{Solución}

Los vectores normales son:

\[
n_1=(5,-14,2),
\]

\[
n_2=(10,-11,2).
\]

Así:

\[
n_1\cdot n_2
=
208.
\]

Además:

\[
\|n_1\|
=
\|n_2\|
=
15.
\]

Por tanto:

\[
\cos\theta
=
\frac{208}{225}.
\]

Así:

\[
\boxed{
\theta
=
\arccos\left(
\frac{208}{225}
\right)
}
\]

\section{Ejercicio 12}
\subsection{Solución}

La recta pasa por:

\[
P_r=(2,3,4)
\]

y tiene vector director:

\[
v=(4,4,2).
\]

El punto dado es:

\[
P=(1,2,3).
\]

La distancia es:

\[
d(P,r)
=
\frac{
\|\overrightarrow{P_rP}\times v\|
}{
\|v\|
}.
\]

Se tiene:

\[
\overrightarrow{P_rP}
=
(-1,-1,-1).
\]

Por tanto:

\[
\boxed{
d(P,r)
=
\frac{\sqrt2}{3}
}
\]

\section{Ejercicio 13}
\subsection{Solución}

La recta es:

\[
\begin{cases}
x=2+3\lambda,\\
y=-1+\lambda,\\
z=1-2\lambda.
\end{cases}
\]

Eliminando \(\lambda\):

\[
x-3y-5=0,
\]

\[
z+2y+1=0.
\]

Por tanto:

\[
\boxed{
r:
\begin{cases}
x-3y-5=0,\\
z+2y+1=0.
\end{cases}
}
\]

Para calcular la distancia al punto:

\[
P=(1,3,-2),
\]

tomamos:

\[
P_r=(2,-1,1),
\qquad
v=(3,1,-2).
\]

Entonces:

\[
d(P,r)
=
\frac{
\|\overrightarrow{P_rP}\times v\|
}{
\|v\|
}.
\]

Se obtiene:

\[
\boxed{
d(P,r)
=
\frac{3\sqrt{10}}{2}
}
\]

\section{Ejercicio 14}
\subsection{Solución}

Sea:

\[
P=(3,1,-2)
\]

y:

\[
\pi:
2x+y-z+1=0.
\]

\textbf{Primera forma: fórmula directa.}

\[
d(P,\pi)
=
\frac{
|2\cdot3+1-(-2)+1|
}{
\sqrt{2^2+1^2+(-1)^2}
}.
\]

Así:

\[
d(P,\pi)
=
\frac{10}{\sqrt6}.
\]

\textbf{Segunda forma: proyección sobre el vector normal.}

Un vector normal es:

\[
n=(2,1,-1).
\]

Tomando un punto \(Q\in\pi\), la distancia es:

\[
d(P,\pi)
=
\left|
\operatorname{proy}_{n}
\overrightarrow{QP}
\right|.
\]

El resultado es de nuevo:

\[
\frac{10}{\sqrt6}.
\]

\textbf{Tercera forma: recta perpendicular.}

La recta perpendicular a \(\pi\) que pasa por \(P\) es:

\[
(x,y,z)
=
(3,1,-2)
+
\lambda(2,1,-1).
\]

Su intersección con \(\pi\) se produce en:

\[
\lambda=-\frac53.
\]

La distancia entre ambos puntos es:

\[
\left|
-\frac53
\right|
\sqrt6.
\]

Por tanto:

\[
\boxed{
d(P,\pi)
=
\frac{5\sqrt6}{3}
}
\]

\section{Ejercicio 15}
\subsection{Solución}

El plano es:

\[
(x,y,z)
=
(0,0,-4)
+
\lambda(2,2,-1)
+
\mu(-3,2,0).
\]

Un vector normal es:

\[
n
=
(2,2,-1)\times(-3,2,0)
=
(2,3,10).
\]

Por tanto, su ecuación general es:

\[
2x+3y+10(z+4)=0.
\]

Así:

\[
\boxed{
2x+3y+10z+40=0
}
\]

Para:

\[
P=(5,5,3),
\]

la distancia es:

\[
d(P,\pi)
=
\frac{
|10+15+30+40|
}{
\sqrt{4+9+100}
}.
\]

Por tanto:

\[
\boxed{
d(P,\pi)
=
\frac{95}{\sqrt{113}}
}
\]

\section{Ejercicio 16}
\subsection{Solución}

Los planos son:

\[
\pi_1:
2x-y-2z+5=0
\]

y:

\[
\pi_2:
4x-2y-4z+15=0.
\]

Dividiendo la segunda ecuación entre \(2\):

\[
\pi_2:
2x-y-2z+\frac{15}{2}=0.
\]

Por tanto:

\[
d(\pi_1,\pi_2)
=
\frac{
\left|
\frac{15}{2}-5
\right|
}{
\sqrt{2^2+(-1)^2+(-2)^2}
}.
\]

Así:

\[
\boxed{
d(\pi_1,\pi_2)
=
\frac56
}
\]

\section{Ejercicio 17}
\subsection{Solución}

Sea:

\[
P=(1,0,0)
\]

y el plano:

\[
x+2y+3z-15=0.
\]

Un vector normal es:

\[
n=(1,2,3).
\]

El pie de la perpendicular es:

\[
H
=
P
-
\frac{
1+2\cdot0+3\cdot0-15
}{
1^2+2^2+3^2
}
n.
\]

Así:

\[
H=(2,2,3).
\]

El punto simétrico es:

\[
P'
=
2H-P.
\]

Por tanto:

\[
\boxed{
P'=(3,4,6)
}
\]

\section{Ejercicio 18}
\subsection{Solución}

La recta \(r\) es:

\[
(x,y,z)
=
(1,-3,0)
+
\lambda(2,-1,1).
\]

La recta \(s\) es:

\[
(x,y,z)
=
(0,1,3)
+
\mu(1,-1,2).
\]

Tomamos:

\[
v_r=(2,-1,1),
\qquad
v_s=(1,-1,2).
\]

Entonces:

\[
v_r\times v_s
=
(-1,-3,-1).
\]

Además:

\[
\overrightarrow{P_rP_s}
=
(-1,4,3).
\]

La distancia entre dos rectas reversas es:

\[
d(r,s)
=
\frac{
\left|
\overrightarrow{P_rP_s}
\cdot
(v_r\times v_s)
\right|
}{
\|v_r\times v_s\|
}.
\]

Así:

\[
\boxed{
d(r,s)
=
\frac{14}{\sqrt{11}}
}
\]

\section{Ejercicio 19}
\subsection{Solución}

La recta \(r\) pasa por:

\[
A=(1,2,-1)
\]

y tiene vector director:

\[
v_r
=
B-A
=
(1,-1,1).
\]

La recta \(s\) pasa por:

\[
C=(4,1,0)
\]

y tiene vector director:

\[
v_s
=
D-C
=
(-4,0,-2).
\]

Sean:

\[
R=A+\lambda v_r,
\]

\[
S=C+\mu v_s.
\]

Para que \(R\) y \(S\) sean los puntos más cercanos:

\[
(R-S)\cdot v_r=0,
\]

\[
(R-S)\cdot v_s=0.
\]

Resolviendo:

\[
\lambda=\frac23,
\qquad
\mu=\frac12.
\]

Por tanto:

\[
R=
\left(
\frac53,
\frac43,
-\frac13
\right),
\]

\[
S=(2,1,-1).
\]

Un vector director de la recta que los une es:

\[
S-R
=
\frac13(1,-1,-2).
\]

Por tanto:

\[
\boxed{
(x,y,z)
=
\left(
\frac53,
\frac43,
-\frac13
\right)
+
t(1,-1,-2)
}
\]

\section{Ejercicio 20}
\subsection{Solución}

La simetría respecto de la recta:

\[
x=y=z
\]

es un giro de \(180^\circ\) alrededor del vector:

\[
u=(1,1,1).
\]

La matriz asociada es:

\[
R
=
\frac23
\begin{pmatrix}
1&1&1\\
1&1&1\\
1&1&1
\end{pmatrix}
-
I.
\]

El plano original es:

\[
x+2y-3z=1.
\]

Su vector normal es:

\[
n=(1,2,-3).
\]

Transformándolo:

\[
Rn=(-1,-2,3).
\]

Por tanto, el plano simétrico es:

\[
-x-2y+3z=1.
\]

Así:

\[
\boxed{
x+2y-3z=-1
}
\]

\section{Ejercicio 21}
\subsection{Solución}

La recta está dada por:

\[
\begin{cases}
x+3y-z=3,\\
2x-y+z=2.
\end{cases}
\]

Los vectores normales son:

\[
n_1=(1,3,-1),
\qquad
n_2=(2,-1,1).
\]

Por tanto, un vector director es:

\[
v
=
n_1\times n_2
=
(2,-3,-7).
\]

El punto:

\[
P=(1,1,1)
\]

pertenece a la recta y también al eje:

\[
x=y=z.
\]

Por tanto, permanece fijo por la simetría.

La simetría del vector director respecto del eje \(x=y=z\) es:

\[
v'
=
(-22,-7,5)
\]

salvo un factor escalar.

Por tanto, la recta simétrica es:

\[
\boxed{
(x,y,z)
=
(1,1,1)
+
\lambda(-22,-7,5)
}
\]

\section{Ejercicio 22}
\subsection{Solución}

El haz es:

\[
(2\lambda+1)x
+
(-3\lambda+2)y
+
(7\lambda-1)
=
0.
\]

Agrupamos los términos con \(\lambda\):

\[
\lambda(2x-3y+7)
+
(x+2y-1)
=
0.
\]

Para que un punto pertenezca a todas las rectas deben cumplirse:

\[
2x-3y+7=0,
\]

\[
x+2y-1=0.
\]

Resolviendo:

\[
x=-\frac{11}{7},
\qquad
y=\frac97.
\]

Por tanto:

\[
\boxed{
P=
\left(
-\frac{11}{7},
\frac97
\right)
}
\]

es el punto fijo del haz.

\section{Ejercicio 23}
\subsection{Solución}

Las rectas son:

\[
r:
mx+(2m-1)y+3=0,
\]

\[
s:
(4m-7)x-(m+2)y-8=0.
\]

Dos rectas son perpendiculares si sus vectores normales son ortogonales tras intercambiar la condición equivalente:

\[
m(4m-7)
+
(2m-1)(-m-2)
=
0.
\]

Simplificando:

\[
2(m^2-5m+1)=0.
\]

Por tanto:

\[
m^2-5m+1=0.
\]

Así:

\[
\boxed{
m
=
\frac{5\pm\sqrt{21}}{2}
}
\]

\section{Ejercicio 24}
\subsection{Solución}

La recta eje de simetría es:

\[
s:
x-y-6=0.
\]

La simetría de un punto \((x,y)\) respecto de \(s\) es:

\[
(x',y')
=
(y+6,x-6).
\]

La recta original es:

\[
r:
2x+4y-7=0.
\]

Como la simetría es involutiva, sustituimos:

\[
x=y'+6,
\qquad
y=x'-6.
\]

Entonces:

\[
2(y'+6)
+
4(x'-6)
-
7
=
0.
\]

Simplificando:

\[
4x'+2y'-19=0.
\]

Por tanto:

\[
\boxed{
4x+2y-19=0
}
\]

\section{Ejercicio 25}
\subsection{Solución}

La ecuación es:

\[
y^2-7xy+4x^2=0.
\]

Representa dos rectas que pasan por el origen:

\[
y=m_1x,
\qquad
y=m_2x,
\]

donde:

\[
m_{1,2}
=
\frac{7\pm\sqrt{33}}{2}.
\]

Las bisectrices forman el mismo ángulo con ambas rectas.

Si:

\[
y=mx,
\]

la condición de igualdad de ángulos conduce a:

\[
7m^2-6m-7=0.
\]

Por tanto:

\[
m
=
\frac{3\pm\sqrt{58}}{7}.
\]

Así, las bisectrices son:

\[
\boxed{
y
=
\frac{3+\sqrt{58}}{7}x
}
\]

y:

\[
\boxed{
y
=
\frac{3-\sqrt{58}}{7}x
}
\]

Equivalentemente, ambas vienen dadas por:

\[
\boxed{
7y^2-6xy-7x^2=0
}
\]

\section{Ejercicio 26}
\subsection{Solución}

Sean:

\[
A=(1,4),
\quad
B=(-2,3),
\quad
C=(5,2).
\]

La mediatriz de \(AB\) y la mediatriz de \(AC\) se cortan en el circuncentro.

Resolviendo sus ecuaciones se obtiene:

\[
\boxed{
O=(1,-1)
}
\]

Comprobamos:

\[
OA^2
=
(1-1)^2+(-1-4)^2
=
25,
\]

\[
OB^2
=
(1+2)^2+(-1-3)^2
=
25,
\]

\[
OC^2
=
(1-5)^2+(-1-2)^2
=
25.
\]

Por tanto:

\[
\boxed{
O=(1,-1)
}
\]

\section{Ejercicio 27}
\subsection{Solución}

Sean:

\[
A=(0,0),
\quad
B=(14,0),
\quad
C=(5,12).
\]

Como \(AB\) es horizontal, la altura desde \(C\) es:

\[
x=5.
\]

La pendiente de \(BC\) es:

\[
m_{BC}
=
\frac{12}{5-14}
=
-\frac43.
\]

Por tanto, la altura desde \(A\) tiene pendiente:

\[
\frac34.
\]

Su ecuación es:

\[
y=\frac34x.
\]

Intersecando con:

\[
x=5,
\]

se obtiene:

\[
y=\frac{15}{4}.
\]

Así:

\[
\boxed{
H=
\left(
5,
\frac{15}{4}
\right)
}
\]

\section{Ejercicio 28}
\subsection{Solución}

El baricentro es la media aritmética de los vértices:

\[
G
=
\left(
\frac{x_A+x_B+x_C}{3},
\frac{y_A+y_B+y_C}{3}
\right).
\]

Por tanto:

\[
G
=
\left(
\frac{3+5+8}{3},
\frac{4+12+15}{3}
\right).
\]

Así:

\[
\boxed{
G
=
\left(
\frac{16}{3},
\frac{31}{3}
\right)
}
\]

\section{Ejercicio 29}
\subsection{Solución}

La diagonal conocida es:

\[
d:
x+2y-5=0.
\]

El vértice dado es:

\[
A=(2,-1).
\]

Como \(A\notin d\), la recta \(d\) es la diagonal que une los otros dos vértices.

El centro del cuadrado es la proyección ortogonal de \(A\) sobre \(d\).

Se obtiene:

\[
O=(3,1).
\]

El vértice opuesto a \(A\) es su simétrico respecto de \(O\):

\[
C=2O-A.
\]

Así:

\[
\boxed{
C=(4,3)
}
\]

La distancia:

\[
AO=\sqrt5.
\]

Un vector director de \(d\) es:

\[
(2,-1),
\]

cuya norma también es:

\[
\sqrt5.
\]

Por tanto, los otros vértices son:

\[
B=O+(2,-1)=(5,0),
\]

\[
D=O-(2,-1)=(1,2).
\]

Así:

\[
\boxed{
B=(5,0),
\qquad
C=(4,3),
\qquad
D=(1,2)
}
\]

La diagonal \(AC\) mide:

\[
2\sqrt5.
\]

El área del cuadrado es:

\[
A
=
\frac{(2\sqrt5)^2}{2}.
\]

Por tanto:

\[
\boxed{
A=10
}
\]

\section{Ejercicio 30}
\subsection{Solución}

Sean:

\[
A=(0,0,0),
\quad
B=(1,0,0),
\quad
C=(0,1,1).
\]

Los vectores de la base son:

\[
\overrightarrow{AB}=(1,0,0),
\]

\[
\overrightarrow{AC}=(0,1,1).
\]

Un vector normal al plano \(ABC\) es:

\[
n
=
\overrightarrow{AB}\times\overrightarrow{AC}
=
(0,-1,1).
\]

Por tanto, el plano de la base es:

\[
\boxed{
z-y=0
}
\]

El plano paralelo que pasa por:

\[
P=(0,10,9)
\]

es:

\[
\boxed{
z-y+1=0
}
\]

El vértice:

\[
D=(x,y,z)
\]

equidista de \(A\), \(B\) y \(C\).

De:

\[
DA=DB
\]

se obtiene:

\[
x=\frac12.
\]

De:

\[
DA=DC
\]

se obtiene:

\[
y+z=1.
\]

Además:

\[
z-y+1=0.
\]

Resolviendo:

\[
\boxed{
D=
\left(
\frac12,
1,
0
\right)
}
\]

El área de la base es:

\[
A_{ABC}
=
\frac12
\|
\overrightarrow{AB}\times\overrightarrow{AC}
\|
=
\frac{\sqrt2}{2}.
\]

La altura es la distancia entre los planos paralelos:

\[
h
=
\frac1{\sqrt2}.
\]

Por tanto:

\[
V
=
\frac13A_{ABC}h.
\]

Así:

\[
\boxed{
V=\frac16
}
\]

El simétrico de \(D\) respecto del plano:

\[
z-y=0
\]

se obtiene intercambiando las coordenadas \(y\) y \(z\).

Por tanto:

\[
\boxed{
D'
=
\left(
\frac12,
0,
1
\right)
}
\]

\end{document}
